\documentclass[11pt,a4paper]{amsart}

\RequirePackage[hyphens]{url}

\RequirePackage[OT1]{fontenc}
\RequirePackage{amsthm,amsmath}
\RequirePackage[numbers]{natbib}
\RequirePackage[colorlinks,citecolor=blue,urlcolor=blue]{hyperref}
\RequirePackage{cleveref}

\usepackage{graphicx}
\usepackage{tikz}
\usepackage{enumitem}
\usepackage[utf8]{inputenc}
\usepackage[english]{babel}
\usepackage{latexsym}
\usepackage{amssymb}
\usepackage{amsmath}
\usepackage{pgfplots}
\usepackage{tikz}
\usepackage{graphicx}
\usepackage[most]{tcolorbox}
\usepackage{nameref}
\usepackage{enumitem}
\usepackage{randomwalk}
\usepackage{geometry}


\usetikzlibrary{arrows,calc}


\usetikzlibrary{automata, positioning}


\tcbuselibrary{theorems}



\newtheorem{theorem}{Theorem}
\newtheorem{lemma}[theorem]{Lemma}
\newtheorem*{lemma*}{Lemma}
\newtheorem{proposition}[theorem]{Proposition}
\newtheorem{corollary}[theorem]{Corollary}
\newtheorem{conjecture}[theorem]{Conjecture}
\newtheorem{definition}[theorem]{Definition}
\newtheorem{remark}[theorem]{Remark}
\newtheorem{assumption}[theorem]{Assumption}
\newtheorem{claim}[theorem]{Claim}
\newtheorem{fact}[theorem]{Fact}
\newtheorem*{fact*}{Fact}

\tcbuselibrary{theorems}
\newtcbtheorem[number within=section]{Definitions}{Määritelmä}{breakable,code={\edef\@currentlabelname{#2}},colframe=green!45!black,fonttitle=\bfseries}{Definitions}
\newtcbtheorem[number within=section]{Examples}{Esimerkki}{breakable,code={\edef\@currentlabelname{#2}},colframe=orange!85!black,fonttitle=\bfseries}{Examples}
\newtcbtheorem[number within=section]{Theorems}{Lause}{breakable,code={\edef\@currentlabelname{#2}},colframe=blue!85!black,fonttitle=\bfseries}{Theorems}
\newtcbtheorem[number within=section]{Propositions}{Propositio}{breakable,code={\edef\@currentlabelname{#2}},colframe=teal!85!black,fonttitle=\bfseries}{Propositions}
\newtcbtheorem[number within=section]{Lemmas}{Lemma}{breakable,code={\edef\@currentlabelname{#2}},colframe=cyan!85!black,fonttitle=\bfseries}{Lemmas}
\newtcbtheorem{Comments}{Huomautus}{breakable,code={\edef\@currentlabelname{#2}},colframe=pink!85!black,fonttitle=\bfseries}{Comments}
\newtcbtheorem{Exercises}{Harjoitustehtävä}{breakable,code={\edef\@currentlabelname{#2}},colframe=yellow!55!black,fonttitle=\bfseries}{Exercises}
\newtcbtheorem[number within=section]{Warnings}{Varoitus}{breakable,code={\edef\@currentlabelname{#2}},colframe=red!65!black,fonttitle=\bfseries}{Warnings}
\newtcolorbox{Proof}{breakable,colframe=gray!55!black,fonttitle=\bfseries,title=Todistus:}
\newtcolorbox{Solution}{breakable,colframe=gray!55!black,fonttitle=\bfseries,title=Ratkaisu:}
\newtcbtheorem[number within=section]{HExercises}{Teoreettinen lisätehtävä}{breakable,code={\edef\@currentlabelname{#2}},colframe=yellow!35!black,fonttitle=\bfseries}{HExercises}

\newtcbtheorem[number within=section]{CExercises}{Laskennallinen lisätehtävä}{breakable,code={\edef\@currentlabelname{#2}},colframe=yellow!75!black,fonttitle=\bfseries}{HExercises}


\usetikzlibrary{decorations.markings}

% axis style, ticks, etc
\pgfplotsset{every axis/.append style={
                    axis x line=middle,    % put the x axis in the middle
                    axis y line=middle,    % put the y axis in the middle
                    axis line style={->}, % arrows on the axis
                    ymajorticks=false,
                    %xlabel={$x$},          % default put x on x-axis
                    %ylabel={$y$},          % default put y on y-axis
                    xtick={-1,1},                     
                    }}
% arrows as stealth fighters
\tikzset{>=stealth}


\newcommand{\T}{\mathcal{T}}
\newcommand{\N}{\mathbb{N}}
\newcommand{\Z}{\mathbb{Z}}
\newcommand{\R}{\mathbb{R}}
\newcommand{\C}{\mathbb{C}}
\newcommand{\E}{\mathbb{E}}
\newcommand{\Q}{\mathbb{Q}}
\newcommand{\F}{\mathcal{F}}
\newcommand{\Tr}{\mathrm{Tr}}
\newcommand{\M}{\mathrm{Markov}}
\newcommand{\1}{\mathbf{1}}
\newcommand{\red}{\color{red}}
\newcommand{\Exp}{\mathrm{Exp}}
\newcommand{\Poi}{\mathrm{Poi}}
\newcommand{\Unif}{\mathrm{Unif}}

\renewcommand{\P}{\mathbb P}
\renewcommand{\S}{\mathcal S}
\renewcommand{\Re}{\mathrm{Re} \,}
\renewcommand{\Im}{\mathrm{Im} \,}
\renewcommand{\epsilon}{\varepsilon}


%\numberwithin{equation}{section}
%\numberwithin{theorem}{section}
%\numberwithin{figure}{section}

\usepackage{eurosym}


\usepackage{fullpage}
%\usepackage[notref, notcite]{showkeys} %Comment to hide all the labels in the file


\addto\captionsenglish{% Replace "english" with the language you use
  \renewcommand{\contentsname}%
    {Sisällysluettelo}%
}

\addto\captionsenglish{% Replace "english" with the language you use
  \renewcommand{\figurename}%
    {Kuva}%
}

\addto\captionsenglish{% Replace "english" with the language you use
  \renewcommand{\abstractname}%
    {Abstrakti}%
}

\addto\captionsenglish{% Replace "english" with the language you use
  \renewcommand{\refname}%
    {Viitteet}%
}

\addto\captionsenglish{% Replace "english" with the language you use
  \renewcommand{\appendixname}%
    {Liite}%
}



	
\setcounter{tocdepth}{1}

\title{Todennäköisyyslaskenta II a, 2022 -- Viikon 1 harjoitustehtävät}

%\author{Christian Webb}

\email{christian.webb@helsinki.fi}


\begin{document}

\maketitle

Tämän viikon harjoituksissa ({\bf 6 tehtävää}) kerrataan hieman joukko-oppia ja tutustutaan kahden ensimmäisen luennon keskeisimpiin käsitteisiin. Ellei tehtävässä muuten mainita,  {\bf kaikki ratkaisut pitää aina perustella}. Perustele päättelysy joko määritelmästä (esim. luentomuistiinpanot), lemmasta tai muusta tuloksesta, tai jostakin ``yleisesti tunnetusta" faktasta, tai sitten jostakin aiemmin päättelemästäsi. {\bf Muista kuitenkin, että pisteitä saa järkevistä yrityksistä (kaiken ei tarvitse olla pilkulleen oikein), ja kannattaa kirjoittaa relevantit määritelmät näkyviin.}

Palauta ratkaisusi Moodlessa pe 16.9. klo 23.59 mennessä (DIGEST työpajaan).

\medskip

Kommenttina tehtävistä: ainakin joillekin opiskelijoille, tehtävät 3, 5 ja 6 saattavat olla raskaimmat. Kannattaa ehkä tehdä muut ensin. {\bf Muista kysyä apua laskareissa ja Telegramissa, jos tehtävät tuntuvat vaikeilta!}


\medskip 

Ensimmäinen tehtävä on silkka lämmittelytehtävä, ja tarkoitus on muistuttaa mieleen relevantit merkinnät ja käsitteet.
\begin{Exercises}{Käsitteiden kertausta}{rev}
Mitkä seuraavista väitteistä pitävät paikkansa (ei tarvitse perustella)? 
\begin{enumerate}
\item Jos perusjoukkomme on $\Omega=\Z=\{0,1,-1,2,-2,...\}$ (kokonaislukujen joukko),  $\{0\}$ on tapahtuma. 

    It is an elementary event.

\item Jos perusjoukkomme on $\Omega=\Z$, tapahtumat $A=\{x\in \Z: x\geq 0\}$ ja $B=\{x\in \Z: x\leq 0\}$ ovat erilliset.

    No, $ 0 \in A$ and $ 0 \in B$.

\item Perusjoukolle $\Omega=\Z$, jono $(p_\omega)_{\omega\in \Z}$, $p_\omega=\frac{1}{1+|\omega|}$ voi olla jonkin jakauman pistemassafunktio.
    Let's assume for contradiction that $ p_\omega$ can be the point mass function of some distribution.
We will use lemma 2.14 from the lecture notes to arrive to a contradiction.
Since $ \Z$ is countable, we can assign every element $ \omega \in \Omega = \Z$ to a natural number
$ i$ and so our basic set $ \Omega = \{ \omega_1, \omega_2, \omega_3,  \dots\} $ has corresponding
probabilities $ \{ p_1, p_2, p_3,\dots \} $ which are equal to 
\begin{equation}
    \label{eq:1}
    p_i = \frac{1}{1 + \omega_i} 
\end{equation}
So far so good. All $ p_i$ are non-negative so the first condition of lemma 2.14 is satisfied.
However as we will see, the second condition, which is the total sum of the probabilities equaling
to 1, will not be satisfied. Assuming that the second condition holds, then 
\begin{align}
    \sum_{i=1}^{\infty} p_i &= \sum_{i=1}^{\infty} \frac{1}{1 + |\omega_i|} \\
			    &= \frac{1}{1+|0|} + \frac{1}{1+|1|} + \frac{1}{1+ |-1|} + \frac{1}{1+|2|} + \frac{1}{1 + |-2|} + \dots \\
			    &= \frac{1}{1 + 0} + 2 \cdot \sum_{i=1}^{\infty} \frac{1}{1+i}. \label{eq:3} 
\end{align}
In the last equation \ref{eq:3} the infinite sum $ \sum_{i=1}^{\infty}$ term diverges, hence 
we get a contradiction to the existence of the sum $ \sum_{i=1}^{\infty} p_i$ meaning that $ p_i$
cannot form the point mass function of some distribution.

\item On olemassa perusjoukkoja $\Omega$‚ tapahtumia $A\subset \Omega$‚ ja todennäköisyysjakaumia $\P$, joille $\P(A)<0$ tai $\P(A)>1$.

    Incorrect, since for a distribution to be a valid probability distribution we need $\P(A) > 0  $. 
Additionally, $ \P(A) \leq 1$ for all $ A \subset \Omega$.

\item Jos perusjoukkomme on jokin epätyhjä joukko $\Omega$, ja meillä on kaksi tapahtumaa $A$ ja $B$, joille $\P(A)=0.1$ ja $\P(B)=0.2$, niin pätee
\[
\P(A\cup B)<0.4.
\]

    According to lemma 3.1 (3), we have that 
\begin{equation}
\P(A \cup B) = \P(A) + \P(B) - \P(A \cap B) = 0.3 - \P(A \cap B).
\end{equation}
Since $ \P(A \cap B) \geq 0$ according to the definition 2.9, we have that $ \P(A \cup B) \leq 0.3 < 0.4$.

\item Jos meillä on äärellisellä (epätyhjällä) perusjoukolla $\Omega$ tasajakauma, niin tapahtuman $A$ todennäköisyys riippuu vain $A$:n alkioiden lukumäärästä ja $\Omega$:n alkioiden lukumäärästä (eikä esimerkiksi siitä mikä joukko $A$ on kyseessä).

    Although the point mass function sequence $ p_i$ would be equal for all  elementary events 
$\omega_i \in \Omega$, the elementary events themselves might not be disjoint. For example, 
for an event $ A = \{ \omega_1, \{\omega_1, \omega_2 \}  $, where $ \omega_1 \cap   \{\omega_1, \omega_2 \} \neq \empty$, 
we have that $ A = \omega_1 \cup  \{\omega_1, \omega_2 \} $ so by lemma 3.1 (3) we see
\begin{align}
    \P(A) = \P(\omega_1 \cup  \{\omega_1, \omega_2 \} ) &= \P(\omega_1) + \P( \{\omega_1, \omega_2
    \} )\\ 
							&- \P(\omega_1 \cap  \{\omega_1, \omega_2 \} ) \\					    
							&= 1/N + 1/N -  \P(\omega_1 \cap \omega_2)
							\\
						       &= 1/N + 1/N - (1/N + 1/N) \\
						       &= 0 
\end{align}
which is not equal to the number of elements inside $ A \subset \Omega$.


\end{enumerate}
{\bf Vihjeitä:} kannattaa vilkaista relevanteja määritelmiä (perusjoukko, tapahtuma, erilliset tapahtumat, pistemassafunktio, todennäköisyysjakauma, tasajakauma), sekä joitakin keskeisiä tuloksia, kuten Propositio 2.11 ja Lause 3.1). Saattaa myös olla hyödyllistä palauttaa mieleen Sarjat-kurssilta harmonisen sarjan käsite, ja mitä siitä tiedetään (sekä sarjojen vertailua).
\end{Exercises}

Jatketaan harjoittelemalla joukko-opillisten operaatioiden käyttämistä tapahtumien rakentamisessa.
\begin{Exercises}{Tapahtumat ja joukko-opilliset operaatiot}{soper}
Olkoon $\Omega$ jokin epätyhjä perusjoukko, ja $A,B,C\subset \Omega$ joitakin tapahtumia. Lausu joukko-opillisten operaatioiden $\cap,\cup,\setminus,\cdot^\mathsf c$ avulla seuraavat tapahtumat (ei tarvitse perustella). 
\begin{enumerate}
\item Kaikki tapahtumat $A$, $B$ ja $C$ tapahtuvat.
\item Vain tapahtuma $A$ tapahtuu (mutta $B$ ja $C$ eivät).
\item Ainakin kaksi tapahtumista $A,B,C$ tapahtuu. 
\end{enumerate}
{\bf Vihje:} Kannattaa kerrata näiden joukko-opillisten operaatioiden määritelmät ja niiden todennäköisyyslaskennalliset tulkinnat. Viimeinen kohta on ehkä tylsä kirjoittaa, mutta silti ihan terveellistä. Kannattaa vain käydä kaikki mahdollisuudet läpi: joko i) $A$ ja $B$ tapahtuvat, tai ii) ... Venn-diagrammit saattavat helpottaa hahmottamista.
\end{Exercises}


Seuraavassa harjoitustehtävässä tutustumme (yksinkertaiseen) satunnaiskävelyyn – tämä on yksi todennäköisyysteorian keskeisimpiä malleja. Sitä käytetään kuvaamaan monia ajasta riippuvia ilmiöitä. Ajatus mallissa on, että meillä on kävelijä, joka heittää kolikkoa. Jos heiton tulos on kruuna, kävelijä ottaa askeleen oikealle, ja jos klaava niin askeleen vasemmalle. Askeleen jälkeen, hän heittää kolikkoa uudelleen, ja jatkaa tätä satunnaista kävelyä. Tämä malli, ja sen variantit, on pohjana hyvin monelle erilaiselle todennäköisyyslaskennan mallille myös sovelluksissa \url{https://en.wikipedia.org/wiki/Random_walk#Applications}. Esimerkiksi yksi finanssimatematiikan yksinkertaisimpia malleja optioiden hinnoittelulle perustuu satunnaiskävelyyn: \url{https://en.wikipedia.org/wiki/Binomial_options_pricing_model}. 

Tässä tehtävässä tutkimme joitakin hyvin yksinkertaisia kysymyksiä malliin liittyen. Vaikka emme vielä ole riippumattomuudesta puhuneet, niin intuition vuoksi, voi pitää mielessä, että alla parametri $q$ kertoo todennäköisyyden, että yksittäisellä hetkellä askel otetaan oikealle (ja $1-q$ on todennäköisyys, että askel otetaan vasemmalle). Tämä tehdään riippumattomasti jokaisessa vaiheessa. Huomaa, että alla jokainen alkeistapahtuma $\omega=(\omega_1,...,\omega_N)$ on jono, joka esittää yksittäistä $N$:n heiton heittosarjan tulosta. $\omega_1$ kuvaa ensimmäisen heiton tulosta, $\omega_2$ toisen heiton, ...  ja $\omega_N$ viimeisen heiton tulosta.


\begin{Exercises}{Satunnaiskävely}{3}
Olkoon $N\geq 1$ jokin positiivinen kokonaisluku ja tarkastellaan perusjoukkoa $\Omega=\{-1,1\}^N=\{(\omega_1,...,\omega_N): \omega_i\in \{-1,1\} \text{ jokaisella } i=1,...,N\}$. Olkoon myös $q\in [0,1]$ ja määritellään kullakin $\omega\in \Omega$
\[
p_\omega=q^{\sum_{i=1}^N\frac{1+ \omega_i}{2}}(1-q)^{\sum_{i=1}^N\frac{1-\omega_i}{2}}.
\]
\begin{enumerate}
\item Osoita, että $(p_\omega)_{\omega\in \Omega}$ toteuttaa Lauseen 2.14 ehdot, eli osoita, että $p_\omega\geq 0$ jokaisella $\omega\in \Omega$ ja $\sum_{\omega\in \Omega}p_\omega=1$.
\item Olkoon $\P$ lukujen $(p_\omega)_{\omega\in \Omega}$ määräämä todennäköisyysjakauma (jonka olemassaolon Lause 2.14 takaa) ja olkoon $A=\{\omega\in \Omega: \omega_1+\omega_2+...+\omega_N=N\}$ (tässä $\omega_1+...+\omega_N$ tulkitaan satunnaiskävelijän sijainniksi $N$:n askeleen jälkeen). Laske  tapahtuman $A$ todennäköisyys $\P(A)$.
\end{enumerate}

{\bf Älä huolestu, jos ensimmäinen kohta tuntuu haastavalta. Muista, että järkevästä yrityksestä saa pisteitä. Myöskin voit tehdä vain toisen kohdan, jos se tuntuu helpommalta.}

\medskip 

{\bf Vihjeitä:} 
\begin{enumerate}
\item Ensimmäisessä kohdassa voi yrittää näyttää (esimerkiksi induktiolla), että  
\[
\sum_{\omega}p_\omega= \left(\sum_{\sigma\in \{-1,1\}}q^{\frac{1+\sigma}{2}}(1-q)^{\frac{1-\sigma}{2}}\right)^N.
\]
Sulkeiden sisällä oleva summassa on kaksi termiä, ja kannattaa yrittää näyttää, että niiden summa on yksi. 

\item Toisessa kohdasssa kannattaa ehkä miettiä, että mitä arvoja lukujen $\omega_1,...,\omega_N\in \{-1,1\}$ tulee saada, jotta ehto $\omega_1+...+\omega_N=N$ toteutuu (on vain yksi $\omega\in \Omega$ jolla tämä toteutuu).
\end{enumerate}
\end{Exercises}


Seuraavan tehtävän on tarkoitus olla hieman todennäköisyysjakauman määritelmän harjoittelua.

\begin{Exercises}{Jakaumien seokset}{4}
Olkoon $\Omega$ jokin (epätyhjä) perusjoukko, ja $\P_1,...,\P_n$ todennäköisyysjakaumia $\Omega$:lla. Olkoon myös $p_1,...,p_n\geq 0$ ja $\sum_{i=1}^n p_i=1$. Osoita, että 
\[
\P(A):=\sum_{i=1}^n p_i \P_i(A)
\] 
määrittelee todennäköisyysjakauman $\Omega$:lla.

\medskip

{\bf Vihje:} Tämän pitäisi olla aika suoraviivaista määritelmän käyttöä.
\end{Exercises}


Seuraavassa harjoitustehtävässä palataan satunnaiskävelyyn. Tämä tehtävä on hieman haastavampi kuin aiemmat. Ei kannata kuitenkaan pelästyä sitä vaan laskea rauhassa loppuun.

\begin{Exercises}{Paluu satunnaiskävelyyn}{5}
Tarkastellaan Harjoituksen \ref{Exercises:3} mallia, ja katsotaan kullakin $k\in \Z$ tapahtumaa 
\[
A=\{\omega\in \Omega: \omega_1+...+\omega_N=0\},
\]
eli, että $N$:n askeleen jälkeen, satunnaiskävelijä päätyy pisteeseen $0$. 

Osoita, että 
\[
\P(A)=\begin{cases}
0, & N \text{ on pariton}\\
{N\choose N/2}q^{\frac{N}{2}}(1-q)^{\frac{N}{2}}, & N \text{ on parillinen}
\end{cases}.
\]

{\bf Vihje:}  

\begin{enumerate}
\item Mieti kuinka monella indeksin arvolla täytyy olla $\omega_j=+1$ (ja kuinka monella $\omega_j=-1$), jotta $\omega_1+...+\omega_N=0$. Tästä pitäisi myös kyetä päättelemään, että miksi parittoman $N$:n tapauksessa kyseinen tapahtuma on mahdoton.
\item Tämän jälkeen, yritä perustella miksi  (parillisella $N$)
\[
\P(A)=q^{\frac{N}{2}}(1-q)^{\frac{N}{2}}|\{\omega\in \Omega:\omega_1+...+\omega_N=0\}|.
\]
\item Lopulta käyttämällä tietoa, että kuinka monella indeksillä $j$ pätee $\omega_j=1$, yritä tulkita yllä olevan joukon alkioiden lukumäärä kysymyksenä, että kuinka monella tavalla voimme poimia joukosta $\{1,...,N\}$ $x$ määrä indeksiä (missä $x$ on niiden indeksien $j$ lukumäärä, joilla $\omega_j=1$). Järjestyksellä ei tässä poiminnassa ole väliä, joten binomikertoimet sopivat tähän erinomaisesti.
\end{enumerate}
\end{Exercises}


Lopuksi tarkastellaan hieman Lauseen 3.1 todennäköisyyslaskennallisten epäyhtälöiden yleistyksiä/laajennuksia. Jälkimmäinen kohta saattaa olla hieman haastava, mutta ei kannata säikähtää sitäkään.
\begin{Exercises}{Todennäköisyyslaskennallisia epäyhtälöitä}{6}
\begin{enumerate}
\item Osoita, että mille tahansa (ei välttämättä erillisille) tapahtumille $(A_j)_{j=1}^n$ pätee
\[
\P\left(\bigcup_{j=1}^n A_j\right)\leq \sum_{j=1}^n \P(A_j).
\]
\item Osoita, että mille tahansa (ei välttämättä erillisille) tapahtumille $(A_j)_{j=1}^\infty$ pätee 
\[
\P\left(\bigcup_{j=1}^\infty A_j\right)\leq \sum_{j=1}^\infty \P(A_j)
\]
\end{enumerate}

{\bf Vihje:} Ensimmäisessä kohdassa kannattaa yrittää induktiota ja Lausetta 3.1, joka antaa tuloksen tapauksessa $n=2$. Toisessa kohdassa voi yrittää lähestyä ongelmaa kahdella tapaa: voi joko yrittää perustella (monotonisuuden avulla) miksi ensimmäisen kohdan tuloksesta voi ottaa $n\to \infty$ rajan, tai sitten voi yrittää kirjoittaa tapahtuman $\bigcup_{j=1}^\infty A_j$ erillisten tapahtumien yhdisteinä (tähän hyviä tapahtumia voisivat olla $B_1=A_1$ ja $B_n=A_n\setminus \bigcup_{j=1}^{n-1}A_j$ kun $n\geq 2$), ja arvioida vastaavia todennäköisyyksiä.
\end{Exercises}



\end{document}
